\documentclass[../main.tex]{subfiles}

\begin{document}

\chapter*{一些记号}

这里给出一些本笔记中的常用记号. 

\begin{enumerate}

\item 对集合 $A$，记 $\mathcal{P}(A)$ 为 $A$ 的幂集，即 $A$ 的所有子集的集合. 

\item 对函数 $f:X \to Y$，集合 $A \subset Y$，$f^{-1}(A) := \{x \in X \mid f(x) \in A\}$ 称为 $A$ 在 $f$ 下的原像. 

\item 符号函数定义为 $\mathrm{sgn}(x) = \begin{cases}
    1, & x > 0; \\
    0, & x = 0; \\
    -1, & x < 0. 
\end{cases}$

\item $\mathbb N$：自然数集，注意 $0 \in \mathbb N$. 

\item $\mathbb Z$：整数环. 

\item $\mathbb Q$：有理数域. 

\item $\mathbb R$：实数域，严格的构造将在 \ref{sec:1.3} 节中给出. 

\item $\mathbb R^n$：$n$ 个 $\mathbb R$ 的笛卡儿积，若不加说明，通常配以 Euclid 内积成为内积空间. 

\item $\overline{\mathbb R}$：广义实数集，即在 $\mathbb R$ 中添加 $-\infty$ 与 $+\infty$ 两个元素得到的集合. 在 $\overline{\mathbb R}$ 上规定：对任意 $x \in \mathbb R$，$x + (\pm\infty) = \pm\infty, x \cdot (\pm\infty) = (\pm\mathrm{sgn}(x))\infty, -\infty < x < +\infty$ 以及 $(\pm\infty) + (\pm\infty) = \pm\infty$，但 $(\pm\infty) - (\pm\infty)$ 与 $0 \cdot (\pm\infty)$ 是不被定义的. 同时，在所涉及的表达式均有定义时，加法与乘法满足与 $\mathbb R$ 上相同的运算律. 

\item $\mathbb C$：复数域. 可认为是在 $\mathbb R^2$ 上定义加法与乘法：$(x_1, y_1) + (x_2, y_2) = (x_1 + x_2, y_1 + y_2), (x_1, y_1) \cdot (x_2, y_2) = (x_1 x_2 - y_1 y_2, x_1 y_2 + x_2 y_1)$，从而使之成为一个域. 将其上的元素 $(x, y)$ 形式地记为 $x + y\i$，且 $\i^2 = -1$. 对 $z = x + y\i \in \mathbb C$，定义其共轭为 $\overline{z} = x - y\i$，同时定义其绝对值（模长）为 $|z| = \sqrt{x^2 + y^2} = \sqrt{z \overline{z}}$. 

\item 对 $a, b \in \overline{\mathbb R}, a < b$，定义开区间 $(a, b) = \{x \in \mathbb R \mid a < x < b\}$，闭区间 $[a, b] = \{x \in \mathbb R \mid a \leqslant x \leqslant b\}$，类似可定义 $(a, b]$ 与 $[a, b)$，规定 $\pm\infty$ 不能作为闭区间的端点，即不允许形如 $[a, +\infty], (a, +\infty]$ 等的区间. 
\end{enumerate}

\end{document}